% GENERATED FILE. Edit canonical lessons, then run npm run generate:books.
% canonical-source-hash: fnv1a64:3ba36a951f33f594
% canonical-lessons: TA-C383-cattiyam, TA-C383-nikaltakavu, TA-C383-anekamaka, TA-C383-cantarppam, TA-C383-illaiyel

\chapter{What Might Happen}
\label{ch:ta-what-might-happen}
\hlchaptermodality{383}

\begin{chapteropening}
\textbf{By the end of this chapter:} \emph{I can say how likely something is.}

Complete the last lesson of chapter 383: I can say how likely something is.
\end{chapteropening}

\section[\texorpdfstring{\ta{சாத்தியம்} (cāttiyam)}{சாத்தியம் (cāttiyam)}]{\ta{சாத்தியம்} (cāttiyam) — a possibility; possible}

\label{lesson:TA-C383-cattiyam}



\begin{quote}
Before the new one: say the Tamil for an exception, then the Tamil for if need be.
\end{quote}

\subsection*{You'll want to know: \ta{சாத்தியம்}}
\textbf{\ta{சாத்தியம்}} — \emph{cāttiyam} — \textquotedblleft{}a possibility; possible\textquotedblright{}.

\begin{quote}
\ta{இந்த} \ta{வேலையை} \ta{ஒரே} \ta{நாளில்} \ta{முடிப்பது} \ta{சாத்தியம்} \ta{இல்லை}.
\end{quote}

\begin{quote}
\emph{inta vēlaiyai orē nāḷil muṭippatu cāttiyam illai.}
\end{quote}

\begin{quote}
It is not possible to finish this work in one day.
\end{quote}

\ta{சாத்தியமான}, possible; \ta{சாத்தியமில்லை}, it is not possible.

\begin{cousinweb}
\textbf{\ta{சாத்தியம்}} (\emph{cāttiyam}) is borrowed from Sanskrit \emph{sādhya}, \textquotedblleft{}to be accomplished; practicable, feasible, attainable\textquotedblright{} (Monier-Williams).
\end{cousinweb}

\begin{grammarlens}[title={What you've built}]
One more word for saying what happens if.
\end{grammarlens}

\subsection*{Your turn}
\begin{itemize}
\raggedright
  \item \emph{Say it:} \emph{cāttiyam}
  \item \emph{Say it:} \emph{cāttiyam}, once more
  \item \emph{Say it:} \emph{inta vēlaiyai orē nāḷil muṭippatu cāttiyam illai.} — It is not possible to finish this work in one day.
  \item \emph{Recall:} say the Tamil for an exception, then the Tamil for if need be, then say \emph{cāttiyam} again
\end{itemize}

\begin{tcolorbox}[breakable,colback=teal!4,colframe=teal!35!black,title={Before you move on}]
What does \ta{சாத்தியம்} mean? (\textquotedblleft{}A possibility; possible\textquotedblright{}.) Say it once more.
\end{tcolorbox}
\section[\texorpdfstring{\ta{நிகழ்தகவு} (nikaḻtakavu)}{நிகழ்தகவு (nikaḻtakavu)}]{\ta{நிகழ்தகவு} (nikaḻtakavu) — probability}

\label{lesson:TA-C383-nikaltakavu}



\begin{quote}
Before the new one: say the Tamil for if need be, then the Tamil for a possibility.
\end{quote}

\subsection*{You'll want to know: \ta{நிகழ்தகவு}}
\textbf{\ta{நிகழ்தகவு}} — \emph{nikaḻtakavu} — \textquotedblleft{}probability\textquotedblright{}.

\begin{quote}
\ta{நாளை} \ta{மழை} \ta{பெய்வதற்கான} \ta{நிகழ்தகவு} \ta{அதிகம்}.
\end{quote}

\begin{quote}
\emph{nāḷai maḻai peyvataṟkāṉa nikaḻtakavu atikam.}
\end{quote}

\begin{quote}
The probability of rain tomorrow is high.
\end{quote}

A word of mathematics; everyday speech says \ta{வாய்ப்பு}, chance.

\begin{cousinweb}
\textbf{\ta{நிகழ்தகவு}} (\emph{nikaḻtakavu}) is a modern native coinage: \emph{nikaḻ}, \textquotedblleft{}to happen, occur\textquotedblright{} (DEDR 3661), and \emph{takavu}, \textquotedblleft{}fitness, proportion\textquotedblright{}, from \emph{taku} (DEDR 3005): the measure of happening.
\end{cousinweb}

\begin{grammarlens}[title={What you've built}]
One more word for saying what happens if.
\end{grammarlens}

\subsection*{Your turn}
\begin{itemize}
\raggedright
  \item \emph{Say it:} \emph{nikaḻtakavu}
  \item \emph{Say it:} \emph{nikaḻtakavu}, once more
  \item \emph{Say it:} \emph{nāḷai maḻai peyvataṟkāṉa nikaḻtakavu atikam.} — The probability of rain tomorrow is high.
  \item \emph{Recall:} say the Tamil for if need be, then the Tamil for a possibility, then say \emph{nikaḻtakavu} again
\end{itemize}

\begin{tcolorbox}[breakable,colback=teal!4,colframe=teal!35!black,title={Before you move on}]
What does \ta{நிகழ்தகவு} mean? (\textquotedblleft{}Probability\textquotedblright{}.) Say it once more.
\end{tcolorbox}
\section[\texorpdfstring{\ta{அநேகமாக} (anēkamāka)}{அநேகமாக (anēkamāka)}]{\ta{அநேகமாக} (anēkamāka) — probably}

\label{lesson:TA-C383-anekamaka}



\begin{quote}
Before the new one: say the Tamil for a possibility, then the Tamil for probability.
\end{quote}

\subsection*{You'll want to know: \ta{அநேகமாக}}
\textbf{\ta{அநேகமாக}} — \emph{anēkamāka} — \textquotedblleft{}probably\textquotedblright{}.

\begin{quote}
\ta{அநேகமாக} \ta{அவர்} \ta{நாளை} \ta{வருவார்}.
\end{quote}

\begin{quote}
\emph{anēkamāka avar nāḷai varuvār.}
\end{quote}

\begin{quote}
He will probably come tomorrow.
\end{quote}

\ta{அநேக}, many + -\ta{ஆக}: in most ways, probably.

\begin{cousinweb}
\textbf{\ta{அநேகமாக}} (\emph{anēkamāka}) is built on Sanskrit \emph{aneka}, \textquotedblleft{}not one, many\textquotedblright{} (Monier-Williams), with the adverb ending -\ta{ஆக}: \textquotedblleft{}in many ways\textquotedblright{}, so \textquotedblleft{}most likely\textquotedblright{}.
\end{cousinweb}

\begin{grammarlens}[title={What you've built}]
One more word for saying what happens if.
\end{grammarlens}

\subsection*{Your turn}
\begin{itemize}
\raggedright
  \item \emph{Say it:} \emph{anēkamāka}
  \item \emph{Say it:} \emph{anēkamāka}, once more
  \item \emph{Say it:} \emph{anēkamāka avar nāḷai varuvār.} — He will probably come tomorrow.
  \item \emph{Recall:} say the Tamil for a possibility, then the Tamil for probability, then say \emph{anēkamāka} again
\end{itemize}

\begin{tcolorbox}[breakable,colback=teal!4,colframe=teal!35!black,title={Before you move on}]
What does \ta{அநேகமாக} mean? (\textquotedblleft{}Probably\textquotedblright{}.) Say it once more.
\end{tcolorbox}
\section[\texorpdfstring{\ta{சந்தர்ப்பம்} (cantarppam)}{சந்தர்ப்பம் (cantarppam)}]{\ta{சந்தர்ப்பம்} (cantarppam) — an occasion; an opportunity}

\label{lesson:TA-C383-cantarppam}



\begin{quote}
Before the new one: say the Tamil for probability, then the Tamil for probably.
\end{quote}

\subsection*{You'll want to know: \ta{சந்தர்ப்பம்}}
\textbf{\ta{சந்தர்ப்பம்}} — \emph{cantarppam} — \textquotedblleft{}an occasion; an opportunity\textquotedblright{}.

\begin{quote}
\ta{சந்தர்ப்பம்} \ta{கிடைத்தால்} \ta{நான்} \ta{அவரிடம்} \ta{பேசுவேன்}.
\end{quote}

\begin{quote}
\emph{cantarppam kiṭaittāl nāṉ avariṭam pēcuvēṉ.}
\end{quote}

\begin{quote}
If I get the chance, I will speak to him.
\end{quote}

\ta{சந்தர்ப்பம்} \ta{கிடை}, to get the chance; \ta{இந்தச்} \ta{சந்தர்ப்பத்தில்}, on this occasion.

\begin{cousinweb}
\textbf{\ta{சந்தர்ப்பம்}} (\emph{cantarppam}) is borrowed from Sanskrit \emph{saṃdarbha}, \textquotedblleft{}stringing or binding together, arranging\textquotedblright{} (Monier-Williams); Tamil uses it for the moment when things come together.
\end{cousinweb}

\begin{grammarlens}[title={What you've built}]
One more word for saying what happens if.
\end{grammarlens}

\subsection*{Your turn}
\begin{itemize}
\raggedright
  \item \emph{Say it:} \emph{cantarppam}
  \item \emph{Say it:} \emph{cantarppam}, once more
  \item \emph{Say it:} \emph{cantarppam kiṭaittāl nāṉ avariṭam pēcuvēṉ.} — If I get the chance, I will speak to him.
  \item \emph{Recall:} say the Tamil for probability, then the Tamil for probably, then say \emph{cantarppam} again
\end{itemize}

\begin{tcolorbox}[breakable,colback=teal!4,colframe=teal!35!black,title={Before you move on}]
What does \ta{சந்தர்ப்பம்} mean? (\textquotedblleft{}An occasion; an opportunity\textquotedblright{}.) Say it once more.
\end{tcolorbox}
\section[\texorpdfstring{\ta{இல்லையேல்} (illaiyēl)}{இல்லையேல் (illaiyēl)}]{\ta{இல்லையேல்} (illaiyēl) — or else; otherwise}

\label{lesson:TA-C383-illaiyel}



\begin{quote}
Before the new one: say the Tamil for probably, then the Tamil for an occasion.
\end{quote}

\subsection*{You'll want to know: \ta{இல்லையேல்}}
\textbf{\ta{இல்லையேல்}} — \emph{illaiyēl} — \textquotedblleft{}or else; otherwise\textquotedblright{}.

\begin{quote}
\ta{நாம்} \ta{சீக்கிரம்} \ta{புறப்பட} \ta{வேண்டும்}, \ta{இல்லையேல்} \ta{பேருந்தைத்} \ta{தவறவிடுவோம்}.
\end{quote}

\begin{quote}
\emph{nām cīkkiram puṟappaṭa vēṇṭum, illaiyēl pēruntait tavaṟaviṭuvōm.}
\end{quote}

\begin{quote}
We must leave early, or else we will miss the bus.
\end{quote}

A formal written word; the everyday one is \ta{இல்லாவிட்டால்}.

\begin{cousinweb}
\textbf{\ta{இல்லையேல்}} (\emph{illaiyēl}) is native: \emph{illai}, \textquotedblleft{}it is not\textquotedblright{} (DEDR 2559), with the old conditional ending -\ta{ஏல்}, \textquotedblleft{}if\textquotedblright{}: \textquotedblleft{}if it is not so\textquotedblright{}.
\end{cousinweb}

\begin{grammarlens}[title={What you've built}]
One more word for saying what happens if.
\end{grammarlens}

\subsection*{Your turn}
\begin{itemize}
\raggedright
  \item \emph{Say it:} \emph{illaiyēl}
  \item \emph{Say it:} \emph{illaiyēl}, once more
  \item \emph{Say it:} \emph{nām cīkkiram puṟappaṭa vēṇṭum, illaiyēl pēruntait tavaṟaviṭuvōm.} — We must leave early, or else we will miss the bus.
  \item \emph{Recall:} say the Tamil for probably, then the Tamil for an occasion, then say \emph{illaiyēl} again
\end{itemize}

\begin{tcolorbox}[breakable,colback=teal!4,colframe=teal!35!black,title={Before you move on}]
What does \ta{இல்லையேல்} mean? (\textquotedblleft{}Or else; otherwise\textquotedblright{}.) Say it once more.
\end{tcolorbox}
